\section{A Taxonomy of Credit Assignment Methods for LLM Agents}\label{sec:taxonomy}

\citet{zhang2026from} organize LLM credit assignment along two axes. The first is \emph{granularity}: token, segment, step/turn, or multi-agent. The second is \emph{methodology}: Monte-Carlo, temporal-difference, LLM-as-critic or model-based, game-theoretic, or information-theoretic. The survey's snapshot covers 47 methods (41 core plus 6 adjacent enablers). Its thesis is that reasoning-side credit assignment is converging on PRMs and critic-free group comparison, while agentic credit assignment keeps producing new mechanisms: hindsight counterfactuals, privileged critics and turn-level MDPs.

We group methods by mechanism into six families, (a)--(f). For a researcher choosing among them, two further questions matter as much as granularity. First, \emph{what must the environment provide}: a resettable state, recurring states, a known gold answer, a per-step verifier, or privileged information? Second, \emph{where does the extra compute go}: more rollouts, critic training, judge calls, or extra forward passes? Each family below closes with a short profile that answers both. Unless stated otherwise, ``gain'' means improvement over GRPO at an equal number of trajectories. Section~\ref{sec:tax-failures} explains why that qualifier matters.

\subsection{(a) Process rewards and implicit dense rewards}\label{sec:tax-prm}

\emph{Idea.} Train or derive a scorer that rates each step, and use its scores as dense rewards. Early PRMs labeled steps by sampling continuations from each step and checking how often they succeed, as in Math-Shepherd~\citep{wang2024math}\unverified, or used human labels~\citep{lightman2023lets}\unverified. Implicit PRMs learn step scores from outcome labels alone. The theoretical basis is that DPO implicitly learns token-level Q-values~\citep{rafailov2024from}\unverified, and PRIME is a representative method~\citep{cui2025process}\unverified.

\emph{The hacking finding.} PURE~\citep{cheng2025stop} traces PRM-induced reward hacking mainly to the standard \emph{summed} return. Under that return the policy seeks out steps the PRM already scores highly. In PURE's experiments, sum-form variants collapsed around training step~25. A \emph{min-form} value, in which a response is worth its weakest step, stayed stable, and adding a little verifiable reward reduced hacking further. A 2026 follow-up reports that min-form only delays hacking~\citep{uprm2026unsupervised}\unverified.

\emph{In agents.} Per-step correctness labels rarely exist for agents. Agent PRMs that work therefore either learn progress from outcomes or use an environment oracle. AgentPRM estimates turn values with TD and GAE~\citep{xi2025agentprm}\unverified. iStar/OPRL trains an implicit PRM online and fuses it with the outcome reward~\citep{istar2025online}. It argues that token-level PRMs learned from outcomes are ``overly fine-grained for agentic learning,'' amplifying variance as trajectories grow. VPR turns symbolic or algorithmic oracles into verifiable turn-level rewards~\citep{vpr2026verifiable}. In all three, the outcome reward stays dominant.

\emph{Profile.} Needs a step verifier, labels, or an outcome-trained scorer, but no state reset. Extra compute: PRM training plus one scoring pass per step. Low variance, but biased by model error and gameable, so it is safe only as shaping on top of a verifiable outcome. Shown to help in math reasoning, and in agents (WebShop, VisualSokoban, ALFWorld-type tasks) when mixed with the outcome reward.

\subsection{(b) Monte-Carlo rollouts and tree/branching estimators}\label{sec:tax-mc}

\emph{Idea.} To estimate the value of an intermediate state, return to it and sample several continuations. Differences between those values give step advantages without a critic. VinePPO replaces PPO's value network with such ``vine'' estimates at intermediate positions~\citep{kazemnejad2024vineppo}\unverified. SPO cuts a chain of thought at chosen points and computes Monte-Carlo advantages per segment, with chain and tree variants~\citep{guo2025segment}\unverified. TreeRL~\citep{hou2025treerl}\unverified{} and TreePO~\citep{li2025treepo}\unverified{} grow on-policy trees and fork at uncertain tokens.

\emph{Agentic trees.} For agents, branching is presented as a \emph{budget saving}, because sibling branches share their prefix. Tree-GRPO makes each tree node a full agent turn~\citep{ji2025tree}. On multi-hop QA with Qwen2.5-3B, it reports beating chain-based methods while using a quarter of the rollout budget, measured in tokens and tool calls. ARPO~\citep{dong2025agentic}\unverified{} and AEPO~\citep{dong2025agenticentropy}\unverified{} branch adaptively where token entropy spikes after tool calls, and AT$^2$PO adds turn-wise credit on an entropy-expanded turn tree~\citep{at2po2026agentic}\unverified. For stateful sandboxes, Branching Policy Optimization (BPO) snapshots the sandbox at chosen branch points and uses sibling continuations as variance-reducing baselines~\citep{bpo2026branching}. It states its key assumption explicitly as \emph{snapshot fidelity}: restoring a snapshot must reproduce the same transition distribution.

\emph{Memory agents.} Memory-R2's LoGo-GRPO~\citep{yan2026memory} belongs in this family. It keeps GRPO's trajectory-level (``global'') advantage over $n$ rollouts. It adds session-level (``local'') advantages from rerollouts that start at a snapshot of the \emph{same} intermediate memory state. Mem-T's MoT-GRPO similarly back-propagates terminal feedback through a tree of memory operations~\citep{memt2026mem}. Section~\ref{sec:memory} reads LoGo-GRPO closely.

\emph{Profile.} Needs a resettable or forkable state. That is trivial for text and for stateless Python calls (replay the prefix), but requires checkpoint infrastructure for sandboxes. Extra compute: about $K$ partial rollouts per estimated point for vines and segments. Trees instead reallocate a fixed budget and can cost less than i.i.d.\ sampling. Unbiased, but high-variance at small $K$; tree siblings are correlated, and entropy-triggered branching concentrates credit on uncertain turns. Shown to help on math reasoning and on search and multi-hop QA agents.

\subsection{(c) Anchor-state grouping}\label{sec:tax-gigpo}

\emph{Idea.} If the same environment state recurs across the $G$ trajectories sampled for one prompt, compare the actions taken at that ``anchor'' state directly. GiGPO adds this inner, step-level group to GRPO's outer trajectory group and uses discounted returns from each occurrence~\citep{feng2025group}. It needs no extra rollouts and no critic. It reports gains of more than 12\% on ALFWorld and more than 9\% on WebShop over GRPO.

\emph{Limitation.} The authors name their reliance on exact state matching as the main limitation. In grid, room or page-state environments, identical observations recur often. In math with a Python interpreter, the state is the full text history plus tool outputs, so exact recurrence after the initial prompt is rare. Anchor groups then shrink to single trajectories and GiGPO reduces to GRPO. This is our inference: we found no report of GiGPO on code-interpreter tool use. Successors relax exact matching to trajectory graphs or state graphs, such as SALT~\citep{salt2026step}\unverified{} and RewardFlow~\citep{rewardflow2026topology}\unverified.

\emph{Profile.} Needs recurring states, no reset and no verifier. Extra compute is close to zero. Exact anchors give an unbiased within-state comparison; similarity-based anchors trade bias for coverage. Shown to help on ALFWorld and WebShop.

\subsection{(d) Turn-level rewards and critics}\label{sec:tax-turn}

\emph{Idea.} Treat the episode as a turn-level MDP in which one action is one turn or tool call. Then give each turn its own reward or value. The family has three sub-types, each with its own assumption.

\begin{itemize}
  \item \textbf{Verifiable or judged turn rewards.} MT-GRPO combines automatic checks for verifiable turns, such as correct tool-call format or successful execution, with an LLM judge for subjective turns~\citep{wei2025reinforcing}\unverified. Execution checks are free, but they measure validity, not progress.
  \item \textbf{Learned turn-level critics.} ArCHer trains an off-policy turn-level critic alongside an on-policy token-level actor~\citep{zhou2024archer}\unverified. Turn-PPO runs PPO with turn-level importance ratios and turn-level advantages~\citep{turnppo2025turn}\unverified. SWEET-RL trains an \emph{asymmetric} critic that sees privileged training-time information, such as the reference solution~\citep{zhou2025sweet}\unverified.
  \item \textbf{Gold-answer-likelihood signals.} These critic-free signals score each turn by how much it raises a model's probability of the known answer. IGPO uses the policy's own probability (information gain)~\citep{igpo2025information}\unverified, and $\Delta$Belief-RL uses changes in the policy's log-probability of the target~\citep{deltabelief2026intrinsic}\unverified. TRACE puts states at tool-call boundaries, scores each with a \emph{frozen} reference model's log-probability of the gold answer, and rewards each call by the TD change in that value, keeping the outcome reward~\citep{trace2026turn}. A third-party summary reports Qwen3-4B rising from 7.2 to 35.6 on BrowseComp-Plus\unverified.
\end{itemize}

\emph{The 2026 cluster.} Turn-level attribution became crowded in 2026. Examples include TCPO~\citep{tcpo2026turn}, verifier-free outcome potentials~\citep{sop2026self}, and community-indexed entries such as leave-one-turn attribution (LAPO), responsibility-attributed rewards (ECHO) and step-wise peer advantages (SPADER)~\citep{thinkwee2026agentsmeetrl}\unverified. The space has no shared benchmark or comparison protocol, so these methods are rarely compared with one another.

\emph{Profile.} Assumptions differ by sub-type: a gold answer at training time (likelihood signals), a privileged channel (asymmetric critics), or a judge (rubrics). None needs a state reset. Extra compute: one frozen-model forward pass per turn boundary, critic training, or judge calls. Likelihood signals can be gamed by turns that restate or leak the answer, and judges carry their own bias. Shown to help on search and deep-research agents, and on ColBench-style collaborative tasks.

\subsection{(e) Hindsight, counterfactual, Shapley and information-theoretic methods}\label{sec:tax-hindsight}

\emph{Idea.} Attribute the outcome to steps after the fact, either by asking a critic or by perturbing steps.

\begin{itemize}
  \item \textbf{LLM hindsight critics.} HCAPO uses the LLM itself as a post-hoc critic that refines step-level Q-values through hindsight reasoning and counterfactual continuations~\citep{tan2026hindsight}. With Qwen2.5-7B-Instruct it reports ALFWorld 91.4\% versus 77.6\% for GRPO, and WebShop success rising from 66.1\% to 73.8\%. A third-party table gives 90.8 for ALFWorld. RICOL has a reflector LLM give verbal feedback on states, and treats the resulting shift in action probabilities as an advantage estimate~\citep{chen2026retrospective}. It is evaluated on BabyAI, a gridworld, not on LLM reasoning.
  \item \textbf{Counterfactual and Shapley.} The Counterfactual Shapley Value defines causal $\phi$-values on a structural causal model and redistributes reward while preserving the optimal policy~\citep{li2026counterfactual}. It is a general-RL paper with no confirmed LLM experiments. SSVPO applies a sequential Shapley value to step credit for language models~\citep{ssvpo2026effective}\unverified. Exact Shapley values are exponential in the number of segments, so practical methods sample permutations or rollouts.
  \item \textbf{Key-token selection.} Applying DAPO gradients only to the 20\% of tokens with the highest entropy (``forking'' tokens) gave $+11.04$ on AIME'25 and $+7.71$ on AIME'24 over full-token DAPO on Qwen3-32B, and was comparable on Qwen3-8B~\citep{wang2025beyond}. These methods decide \emph{which tokens receive gradient}, not how much credit a step deserves, so they can be combined with turn-level methods.
\end{itemize}

\emph{Profile.} Hindsight critics need an LLM that can judge steps (one extra pass per trajectory); they are not grounded by a verifier and carry judge bias. Counterfactual and leave-one-out methods need resets and cost about as much as Monte-Carlo rollouts. Shapley methods are principled but expensive. Entropy-based token selection is nearly free but ignores errors made at low-entropy tokens. Shown to help on ALFWorld, WebShop, BabyAI and AIME.

\subsection{(f) Learned value models revived}\label{sec:tax-value}

\emph{Idea.} Make the token-level critic work again. VC-PPO attributes PPO's collapse on long chains of thought to value-initialization bias and reward-signal decay~\citep{yuan2025whats}. Its fixes are value pretraining and \emph{decoupled} GAE, with $\lambda=1$ for the critic. VAPO adds length-adaptive GAE, which sets $\lambda$ by sequence length, and reaches 60.4 on AIME 2024 with a Qwen2.5-32B base, beating DAPO~\citep{yue2025vapo}. In agents, learned critics survive mainly at the \emph{turn} level (ArCHer, Turn-PPO, AgentPRM) or as privileged critics (SWEET-RL), where there are far fewer decisions than tokens and the critic's unit matches the unit of decision. We found no controlled evidence that token-level critic PPO beats GRPO-family methods on multi-turn tool use at equal compute.

\emph{Profile.} Needs no reset and no verifier. Extra compute: a critic model plus pretraining. Low variance if tuned, but biased by value error. Shown to help on long-CoT math.

\subsection{Failure modes and audits}\label{sec:tax-failures}

\paragraph{Behavioral hacking around tools}
Dense credit creates incentives for particular tool behaviors. Proof-of-Use documents \emph{tool-call hacking}, in which deep-research agents are rewarded for \emph{invoking} tools rather than using their outputs~\citep{proofofuse2025proof}. A practitioner guide warns that per-call shaping is exploited ``by spamming retrieval calls'' and recommends that the outcome reward dominate~\citep{fireworks2025best}. Aggregator summaries report the opposite failure, indiscriminate tool-call suppression under coarse shaping\unverified. The general rule is that a positive per-call reward invites over-calling and a per-call cost invites under-calling. The Reward Hacking Benchmark associates RL post-training with higher hacking rates in tool-use agents, and simple environment hardening cut hacking by 5.7\,pp without hurting task success~\citep{rewardhacking2026benchmark}.

\paragraph{Is the credit even right?}
Several 2026 efforts test credit signals directly. \emph{Credit Without Ground Truth}~\citep{zhang2026credit} audits step-level credit against \emph{executed replay}, \ie re-running alternative actions to measure their causal effect (ALFWorld, Qwen2.5-7B-Instruct, 50 trajectories, about 20.5K replay rollouts). LLM-judge scores, outcome-conditioned log-probability ratios and policy confidence were no more reliable than a marginal-matched shuffled control, and only 30.5\% of decision points had a measurable causal effect. In a pre-registered seven-arm training study, apparent differences between credit rules reflected \emph{training dose}: sparse credit discards samples, so optimizer step counts differed by an order of magnitude. The authors therefore argue for comparisons at matched effective sample size\unverified. VICT traces credit with verifier instrumentation~\citep{vict2026verifier}\unverified. AgentCreditBench checks turn-level credit from GRPO, RLOO, GAE and GiGPO against exact policy-advantage oracles in small finite-horizon MDPs~\citep{hectopascal2026agentcreditbench}\unverified.

\paragraph{The evaluation protocol}
Most papers in families (a)--(f) report gains at an \emph{equal number of trajectories}. Some methods spend extra rollouts (vines, counterfactuals), some save them (trees), and some discard samples (filtering, sparse credit). At equal trajectories, a method's gain can therefore come from compute or dose rather than from better credit. Tree-GRPO is one of the few that reports equal-token and equal-tool-call budgets~\citep{ji2025tree}. Head-to-head comparisons between credit methods are also rare. According to a citing paper, HCAPO's margin over GiGPO on ALFWorld (about 3--4 points) is much smaller than its margin over GRPO~\citep{vict2026verifier}\unverified. We recommend reporting results at matched trajectories, matched tokens and tool calls, and matched effective samples.
